\documentclass[11pt,twoside,openany]{book}

% --- 1. PAQUETES ---
\usepackage[utf8]{inputenc}
\usepackage[spanish]{babel}
\usepackage[T1]{fontenc}
\usepackage{amsmath, amssymb} 
\usepackage{geometry}
\geometry{top=2cm, bottom=2cm, left=1.5cm, right=1.5cm} % Márgenes para centrado

\usepackage[table]{xcolor}
\usepackage{tabularx}
\usepackage{multicol}
\usepackage{enumitem} 
\usepackage{tcolorbox}
\tcbuselibrary{skins, breakable, raster, listings}

% --- 2. COLORES ---
\definecolor{VoidBlack}{HTML}{0D0D0D}
\definecolor{VoidPurple}{HTML}{3D0075}
\definecolor{BloodPurple}{HTML}{660033}

% --- 3. COMANDO STATS (CENTRADOS Y A LO ANCHO) ---
\newcommand{\stats}[6]{%
    \noindent
    % Usamos >{\centering\arraybackslash} para que los números queden en el centro de la celda
    \begin{tabularx}{\linewidth}{|>{\centering\arraybackslash}X|>{\centering\arraybackslash}X|>{\centering\arraybackslash}X|>{\centering\arraybackslash}X|>{\centering\arraybackslash}X|>{\centering\arraybackslash}X|}
    \hline
    \rowcolor{VoidPurple!30} \textbf{MOV} & \textbf{MELEE} & \textbf{RANGO} & \textbf{REF} & \textbf{BLOQ} & \textbf{FOCUS} \\ \hline
    #1 & #2 & #3 & #4 & #5 & #6 \\ \hline
    \end{tabularx}%
}

% --- 4. FICHA FRENTE (REDISEÑO CON DAÑO LOCALIZADO) ---
\newcommand{\fichaFrente}[9]{%
    \noindent
    \begin{tcolorbox}[enhanced, width=\textwidth, height=12.5cm, colback=white, colframe=VoidBlack, sharp corners, 
        title={\Large \textbf{#1} \hfill \small Coste: #2 pts | #3},
        boxed title style={colback=VoidPurple}]
        
        % 1. ESTADÍSTICAS ARRIBA
        \stats{#4}{#5}{#6}{#7}{#8}{#9}
        
        \vspace{0.3cm} % Separación
        
        % 2. BLOQUE DIVIDIDO (Imagen izq. | Vida y Pasiva der.)
        \begin{minipage}[t]{0.35\linewidth} % Hice la imagen más angosta para dar espacio a la tabla
            \vspace{0pt} 
            \begin{tcolorbox}[enhanced, height=5.5cm, colback=gray!5, colframe=VoidPurple!30, 
                title=REGISTRO VISUAL, boxed title style={colback=VoidBlack, size=small}, nobeforeafter]
                \centering \vfill \textit{[Imagen]} \vfill
            \end{tcolorbox}
        \end{minipage}%
        \hfill
        \begin{minipage}[t]{0.62\linewidth} % Más espacio para la tabla de vida
            \vspace{0pt} 
            
            % NUEVA TABLA DE LOCALIZACIÓN BISKRORUM
            \begin{tcolorbox}[colback=white, colframe=BloodPurple, title=SALUD Y LOCALIZACIÓN (1d100), size=small, nobeforeafter]
                \centering
                \footnotesize
                \begin{tabularx}{\linewidth}{|l|c|X|c|}
                \hline
                \rowcolor{BloodPurple!10} \textbf{ZONA} & \textbf{1d100} & \textbf{PV (VIDA)} & \textbf{RD} \\ \hline
                \textbf{Cuerpo} & 01 - 65 & $\square\square\square\square\square\square$ & -- \\ \hline
                \textbf{Cerebro} & 66 - 100 & $\square\square\square\square$ & -- \\ \hline
                \end{tabularx}
            \end{tcolorbox}
            
            \vspace{0.2cm}
            \begin{tcolorbox}[colback=white, colframe=VoidPurple, title=PASIVA ÚNICA, size=small, nobeforeafter]
                \footnotesize \textbf{Gen\_#1:} Ver detalle al dorso $\rightarrow$
            \end{tcolorbox}
        \end{minipage}

        \vfill % Empuja el equipamiento hacia el fondo

        % 3. EQUIPAMIENTO ABAJO
        \begin{tcolorbox}[colback=gray!5, colframe=VoidBlack, title=MODIFICACIONES GENÉTICAS / ITEMS, size=small, nobeforeafter]
            \begin{tabularx}{\linewidth}{X|c|c}
            \textbf{Gen o Item Equipado} & \textbf{Efecto} & \textbf{Costo} \\ \hline
            1. Arma Base \dotfill & \dotfill & -- \\ \hline
            2. \dotfill & \dotfill & \dotfill \\
            \end{tabularx}
        \end{tcolorbox}
    \end{tcolorbox}%
}

% --- 5. FICHA DORSO (MILIMÉTRICA) ---
\newcommand{\fichaDorso}[3]{%
    \noindent
    \begin{tcolorbox}[enhanced, width=\textwidth, height=12.5cm, colback=white, colframe=VoidBlack, sharp corners, 
        title={\Large #1 \hfill \small RESUMEN DE DATOS},
        boxed title style={colback=VoidBlack}]
        %
        \begin{multicols}{2}
            \section*{Poderes}
            \begin{itemize}[leftmargin=*, noitemsep]
                #2
            \end{itemize}
            
            \columnbreak
            
            \section*{Lore}
            \footnotesize \textit{#3}
        \end{multicols}
        
        \vfill
        \begin{center}
            \footnotesize \textit{Recorta por el borde negro. Esta cara queda detrás del frente al imprimir.}
        \end{center}
    \end{tcolorbox}%
}

\begin{document}
    \chapter*{Codex: Biskrorum}
\begin{center}
    {\Huge \textbf{EL ENJAMBRE ABISAL}} \\
    \vspace{0.5cm}
    {\large \textit{''La Mente es Una. Mis hijos, infinitos.''}}
\end{center}

\section*{El Padre de las Profundidades}
En lo más profundo de Keitros, donde la presión del agua aplastaría el acero, habita una entidad de un plano externo. No es un dios en el sentido convencional; es una conciencia transdimensional que manipula la realidad como si fuera código. Los Biskrorum son sus neuronas externas. No existe el "yo" en su sociedad, solo el "nosotros".

\section*{Naves de Cristal Oscuro}
Cuando el Padre fija su mirada en la superficie, el mar vomita naves de geometría imposible. Para un observador externo, parecen aviones de guerra; para el Padre, son simplemente extremidades con las que golpear. Cada piloto es un hijo que entrega su vida útil, que apenas dura unas horas, para canalizar rayos que desintegran la materia.

\begin{tcolorbox}[colback=VoidPurple!5, colframe=VoidPurple, title=PASIVA DE FACCIÓN: ENLACE NEURONAL]
\textbf{Sincronía del Vacío:} Todas las unidades Biskrorum a 15cm de un Mago pueden repetir sus tiradas de Impacto cuando son Críticos (sacar un 12 natural en la tirada de exito) una vez por ronda.
\end{tcolorbox}
    \newpage
    % --- UNIDAD 1: CAZA SOMBRA-KIEMÖN ---
\newpage % Página Impar (Frente)
\fichaFrente{Caza Sombra-Kiemön}{12}{Nave Ligera}{40cm}{+1}{+5}{+4}{+1}{+6}

\newpage % Página Par (Dorso - Justo detrás)
\fichaDorso{Caza Sombra-Kiemön}{
    \item \textbf{Aceleración:} +2 REF si movió >20cm.
    \item \textbf{Rayo Vacío:} Ignora RD.
    \item \textbf{Suicidio:} Explota al morir (1d6).
}{
    "No es una construcción de metal, sino de voluntad pura. El Sombra-Kiemön es el primer suspiro del Padre al emerger."
}

% --- UNIDAD 2: NÓDULOS DE CHOQUE ---
\newpage % Página Impar (Frente)
\fichaFrente{Nódulos de Choque}{5}{Enjambre}{15cm}{+3}{+2}{+2}{+2}{+3}

\newpage % Página Par (Dorso)
\fichaDorso{Nódulos de Choque}{
    \item \textbf{Biomasa:} Explota al ser retirado.
    \item \textbf{Manto:} Sacrifica modelo por +2 BLOQ.
}{
    "Balas vivientes. Nacen de huevos abisales y regresan al Vacío cuando su biomasa se agota."
}
    \newpage
    \chapter{Mercado Genético del Vacío}

\section*{Evolución por Sacrificio}
Los Biskrorum son débiles en el plano físico (Resistencia Física baja), pero su Resistencia Mágica es abrumadora. Para alcanzar su máximo potencial, deben consumir su propia biomasa. Puedes invertir de \textbf{1 a 3 puntos} por Gen al armar tu lista.

\vspace{0.5cm}

\begin{tcolorbox}[enhanced, colback=white, colframe=VoidBlack, title=CATÁLOGO DE MUTACIONES BÁSICAS]
\begin{tabularx}{\linewidth}{l|c|X}
\rowcolor{BloodPurple!20} \textbf{Gen / Mutación} & \textbf{Pts} & \textbf{Efecto y Sacrificio} \\ \hline

\textbf{Glándula de Furia Abisal} & 3 & \textbf{Sacrificio:} Al inicio de la partida, la unidad pierde 3 PV. \textbf{Efecto:} Gana \textbf{+2 al modificador de Daño} por el resto de la partida. \\ \hline

\textbf{Mutación Óptica} & 2 & \textbf{Sacrificio:} Pierde 3 PV. \textbf{Efecto:} Gana \textbf{+1 al modificador de RANGO} permanente. \\ \hline

\textbf{Espasmo Nervioso} & 1 & \textbf{Activa:} Una vez por turno, puedes sacrificar 1 PV para ganar \textbf{+5cm de MOVIMIENTO} instantáneo. \\ \hline

\textbf{Corteza Refractaria} & 2 &\textbf{Sacrificio:} Al inicio de la partida, la unidad pierde 2 PV. \textbf{Efecto:} Gana \textbf{+2 al modificador de Esquiva} por el resto de la partida. \\ \hline
\end{tabularx}
\end{tcolorbox}
    \newpage
    \cleardoublepage 
\chapter{Grimorio de Combate: Spells y Auras}

\textit{Cartas formato TCG (88x63mm). Diseño optimizado: el cuadro de Costo ancla la esquina inferior izquierda, compartiendo altura con el bloque de Efecto y Lore.}

\vspace{0.5cm}

\begin{center}
\begin{tcbitemize}[
    raster columns=2, 
    raster rows=4,
    raster column skip=0.4cm,
    raster row skip=0.3cm,
    enhanced,
    width=8.8cm, 
    height=6.3cm,
    sharp corners,
    colframe=VoidBlack,
    colback=VoidPurple!5,
    boxed title style={colback=VoidPurple},
    fontupper=\footnotesize,
    halign upper=flush left,
    valign=top, 
    raster width=18cm, 
    raster force size,
    boxsep=2pt, left=2pt, right=2pt, top=0pt, bottom=2pt % bottom=2pt pega el contenido al fondo
]

    % ==========================================
    % CARTA 1: BARRERA ENTRÓPICA
    % ==========================================
    \tcbitem[title={\textbf{BARRERA ENTRÓPICA} \hfill \small 1 PE}]
        
        % 1. Subtítulo
        {\fontsize{6.5}{7}\selectfont \textbf{Tipo:} Reacción \hfill \textbf{Unidad:} Cualquiera} 
        \vspace{2pt}

        % 2. Imagen Placeholder (Ahora es más grande y protagonista)
        \begin{tcolorbox}[colback=gray!15, colframe=gray!40, height=2.6cm, 
                          halign=center, valign=center, sharp corners, size=small, nobeforeafter]
            \textcolor{gray!80}{\textit{[Imagen Placeholder]}}
        \end{tcolorbox}
        
        \vfill % Empuja el bloque inferior hacia la esquina de abajo

        % 3. BLOQUE INFERIOR (Costo a la Izq | Efecto + Leyenda a la Der)
        \noindent
        \begin{minipage}[b]{0.22\linewidth}
            % Cuadrado de Costo (Altura total de 1.8cm)
            \begin{tcolorbox}[colback=BloodPurple, colframe=VoidBlack, sharp corners, 
                              halign=center, valign=center, boxsep=0pt, left=0pt, right=0pt, top=0pt, bottom=0pt, nobeforeafter, height=1.8cm]
                \textcolor{white}{\fontsize{5.5}{6}\selectfont COSTO\\[3pt] \Large \textbf{1}}
            \end{tcolorbox}
        \end{minipage}%
        \hfill
        \begin{minipage}[b]{0.76\linewidth}
            % Caja de Efecto (Altura de 1.3cm)
            \begin{tcolorbox}[colback=white, colframe=VoidPurple, sharp corners, boxrule=0.5pt, 
                              left=2pt, right=2pt, top=2pt, bottom=2pt, nobeforeafter, height=1.3cm]
                \fontsize{6}{7}\selectfont \textbf{Efecto:} Absorbe todo el daño del ataque enemigo. Sufres \textbf{1 Daño Directo} inevitable.
            \end{tcolorbox}
            \vspace{2pt}
            % Leyenda (Ocupa los 0.5cm restantes para igualar al cuadrado)
            \begin{center}
                \fontsize{5.5}{6}\selectfont \textit{''Mejor sangrar por el Padre que por el enemigo.''}
            \end{center}
        \end{minipage}


    % ==========================================
    % CARTA 2: BALIZA DE ENJAMBRE
    % ==========================================
    \tcbitem[title={\textbf{BALIZA DE ENJAMBRE} \hfill \small 2 PE}]
        
        {\fontsize{6.5}{7}\selectfont \textbf{Tipo:} Artefacto \hfill \textbf{Unidad:} Mago} 
        \vspace{2pt}

        \begin{tcolorbox}[colback=gray!15, colframe=gray!40, height=2.6cm, 
                          halign=center, valign=center, sharp corners, size=small, nobeforeafter]
            \textcolor{gray!80}{\textit{[Imagen Placeholder]}}
        \end{tcolorbox}

        \vfill

        \noindent
        \begin{minipage}[b]{0.22\linewidth}
            \begin{tcolorbox}[colback=BloodPurple, colframe=VoidBlack, sharp corners, 
                              halign=center, valign=center, boxsep=0pt, left=0pt, right=0pt, top=0pt, bottom=0pt, nobeforeafter, height=1.8cm]
                \textcolor{white}{\fontsize{5.5}{6}\selectfont COSTO\\[3pt] \Large \textbf{3}}
            \end{tcolorbox}
        \end{minipage}%
        \hfill
        \begin{minipage}[b]{0.76\linewidth}
            \begin{tcolorbox}[colback=white, colframe=VoidPurple, sharp corners, boxrule=0.5pt, 
                              left=2pt, right=2pt, top=2pt, bottom=2pt, nobeforeafter, height=1.3cm]
                \fontsize{6}{7}\selectfont \textbf{Efecto:} Aura 20cm. Otorga Daño Explosivo. Mantenimiento: \textbf{1 FOCUS y 1 PV} por turno.
            \end{tcolorbox}
            \vspace{2pt}
            \begin{center}
                \fontsize{5.5}{6}\selectfont \textit{''Extendiendo la red neuronal del Padre.''}
            \end{center}
        \end{minipage}


    % ==========================================
    % CARTA 3: MANTO OSCURO
    % ==========================================
    \tcbitem[title={\textbf{MANTO OSCURO} \hfill \small 0 PE}]
        
        {\fontsize{6.5}{7}\selectfont \textbf{Tipo:} Aura Soporte \hfill \textbf{Unidad:} Héroes} 
        \vspace{2pt}

        \begin{tcolorbox}[colback=gray!15, colframe=gray!40, height=2.6cm, 
                          halign=center, valign=center, sharp corners, size=small, nobeforeafter]
            \textcolor{gray!80}{\textit{[Imagen Placeholder]}}
        \end{tcolorbox}
        
        \vfill

        \noindent
        \begin{minipage}[b]{0.22\linewidth}
            \begin{tcolorbox}[colback=BloodPurple, colframe=VoidBlack, sharp corners, 
                              halign=center, valign=center, boxsep=0pt, left=0pt, right=0pt, top=0pt, bottom=0pt, nobeforeafter, height=1.8cm]
                \textcolor{white}{\fontsize{5.5}{6}\selectfont COSTO\\[3pt] \Large \textbf{2}}
            \end{tcolorbox}
        \end{minipage}%
        \hfill
        \begin{minipage}[b]{0.76\linewidth}
            \begin{tcolorbox}[colback=white, colframe=VoidPurple, sharp corners, boxrule=0.5pt, 
                              left=2pt, right=2pt, top=2pt, bottom=2pt, nobeforeafter, height=1.3cm]
                \fontsize{6}{7}\selectfont \textbf{Efecto:} Aliados a 10cm ganan +5cm de Movimiento por cada PA. Mantenimiento: \textbf{1 FOCUS}.
            \end{tcolorbox}
            \vspace{2pt}
            \begin{center}
                \fontsize{5.5}{6}\selectfont \textit{''Viajamos por las sombras entre los átomos.''}
            \end{center}
        \end{minipage}

\end{tcbitemize}
\end{center}

\newpage
\thispagestyle{empty}
\mbox{} 
\newpage
\end{document}